% 04 · K-means — the default tool, and its assumptions
\section{K-MEANS: THE DEFAULT TOOL}
\label{sec:kmeans}

\deckslides{slides 4--11 --- the K-means tour, ending on its limits.}

K-means is where every clustering analysis starts, and it is worth
understanding precisely, because its failure modes are the ones the rest of
the toolbox was invented to repair. The method was proposed by
\citet{MacQueen:67} and its modern formulation --- the one everyone implements
--- is due to \citet{Lloyd:82}.

\subsection{The objective}

Given $n$ points $\mathbf{x}_1,\dots,\mathbf{x}_n$ in $\mathbb{R}^d$ and a
fixed number of groups $K$, K-means looks for $K$ centres
$\boldsymbol{\mu}_1,\dots,\boldsymbol{\mu}_K$ and an assignment of every point
to one centre that minimises the \emph{sum of squared errors} (SSE), also
called the within-cluster sum of squares:

\begin{equation}
J = \sum_{k=1}^{K} \sum_{\mathbf{x} \in C_k}
\left\| \mathbf{x} - \boldsymbol{\mu}_k \right\|_2^2 ,
\qquad
\boldsymbol{\mu}_k = \frac{1}{|C_k|}\sum_{\mathbf{x} \in C_k} \mathbf{x}.
\label{eq:sse}
\end{equation}

The second half of the equation is not a separate model choice but a
consequence of the first: for a fixed assignment, the centre that minimises
the sum of squared distances to the members is the arithmetic mean. Two
properties follow immediately and both are assumptions in disguise.

\begin{itemize}
\item Because $J$ penalises distance \emph{squared}, the objective is
  dominated by the points furthest from their centre. A single outlier can
  drag a centre across the field. K-means is not robust; medians (K-medians)
  or medoids are the standard remedy when outliers are expected.
\item Because every point is assigned to its nearest centre, the partition is
  a \emph{Voronoi tessellation}: cells are convex and their boundaries are
  hyperplanes. A cluster with a crescent or dumbbell shape cannot be
  represented, however obvious it looks to the eye
  (\textbf{Figure~\ref{fig:kmeansfail}}).
\end{itemize}

\begin{marginnote}[]
\entry{SSE / inertia}{The sum of squared distances from each point to its
assigned centre: the quantity K-means minimises. Also called inertia in
scikit-learn and within-cluster sum of squares (WCSS).}
\entry{Voronoi cell}{The region of space closer to one centre than to any
other. K-means partitions the plane into Voronoi cells by construction.}
\end{marginnote}

\subsection{The algorithm, and why it always stops}

Lloyd's algorithm is a two-move loop. \emph{Assign} every point to its nearest
centre; \emph{update} every centre to the mean of its members; repeat. Each
move can only reduce $J$: the assignment step is optimal for fixed centres,
and the update step is optimal for fixed assignment. $J$ is bounded below by
zero and takes finitely many values (there are finitely many partitions of
$n$ points into $K$ groups), so the loop must terminate --- in practice after
a few dozen iterations. It terminates at a \emph{local} minimum, and the
distinction matters: K-means is a coordinate-descent algorithm on a
non-convex objective, so different starting centres end in different answers.

\begin{figure}[htbp]\centering
  \fig[0.94]{kmeans_film.png}
  \caption{The K-means loop on a synthetic 2-D dataset, sampled at four
  points in the animation: initial centres, first assignment, first update,
  converged. The workbook prints the key frames; the deck plays the full
  sequence (\textbf{Figure~\ref{fig:qrcodes}}).}
  \label{fig:kmeansfilm}
\end{figure}

\begin{figure}[htbp]\centering
  \begin{tabular}{@{}c@{\hskip6pt}c@{}}
    \fig[0.44]{kmeans_step1.png} &
    \fig[0.44]{kmeans_step2.png} \\[3pt]
    \fig[0.44]{kmeans_step3.png} &
    \fig[0.44]{kmeans_step4.png}
  \end{tabular}
  \caption{The same four states, in the deck's step figures. Stars are the
  centres, dots the data: \panel{a} all points unassigned with $K$ centres
  dropped at random; \panel{b} assignment to the nearest centre; \panel{c}
  centres moved to the mean of their members; \panel{d} the fixed point of the
  iteration --- no point changes hands, so the loop stops.}
  \label{fig:kmeanssteps}
\end{figure}

\subsection{Choosing $K$, and choosing where to start}

$K$ is an input, which is uncomfortable in a problem where the number of
clusters is exactly what we are trying to learn. Three standard answers, and a
warning:

\begin{enumerate}
\item \textbf{Elbow.} Plot $J$ against $K$; $J$ decreases monotonically
  (more centres always fit better), and the curve is expected to bend where
  extra groups stop buying real structure. It is a heuristic: on smooth data
  there is no elbow, and on clustered data with a broad size distribution the
  bend is shallow.
\item \textbf{Silhouette.} For each point, compare its mean distance to its own
  cluster with its mean distance to the nearest other cluster; average over
  points and pick the $K$ that maximises the score \citep{Rousseeuw:87}. It is
  a ratio of within- to between-cluster distance, so it also penalises
  non-convex clusters.
\item \textbf{Model selection.} Compare $K$ by held-out likelihood under a
  Gaussian mixture fitted by EM \citep{Dempster:77}. K-means is the hard
  assignment limit of that model with spherical, equal-variance components ---
  which is precisely the assumption list of \textbf{Table~\ref{tab:assumptions}}.
\end{enumerate}

The warning is that in this project we \emph{cheat}, deliberately. The
catalogue tells us there are 25 clusters, so the benchmarking arms are handed
the true number of groups. That makes the comparison between methods clean,
and it flatters all of them relative to a blind search --- which is why
\S\ref{sec:benchmark} reports field retrieval (a discovery-like task) and
cluster-only separation (a $K$-known task) separately.

Initialisation has the same status: a modelling decision that changes the
answer. The classical approach is to run the loop from several random starts
and keep the best $J$; the modern default is \emph{k-means++}
\citep{Arthur:07}, which samples the first centre uniformly and each
subsequent centre with probability proportional to $D(\mathbf{x})^2$, the
squared distance to the nearest centre already chosen. The effect is to spread
the initial centres out, which typically reaches a better local optimum in
fewer restarts.

\begin{figure}[htbp]\centering
  \fig[0.58]{kmeans_init.png}
  \caption{Initialisation sensitivity (from the deck): the same data, two
  random starts, two different final partitions --- and one of them splits a
  real group. Nothing about the algorithm changed; only the seed did. In a
  clustering paper this is the first thing a referee will ask about, and it is
  why every number in this workbook is reported as a mean over seeds (seven,
  in the head-to-head of \S\ref{sec:spectral}).}
  \label{fig:kmeansinit}
\end{figure}

\subsection{What it assumes, and when it breaks}

\begin{issues}[THE ASSUMPTION LIST]
K-means assumes that clusters are (i) separable by a centre, (ii) roughly
convex, (iii) of comparable size and variance, (iv) all present in equal
proportion, and (v) that $K$ is known. It additionally assumes that the metric
is meaningful --- that the dimensions have been standardised, or that the
scales carry physical meaning. When it fails, it fails in a characteristic
way: it splits large or elongated groups and merges neighbouring small ones,
and it never returns ``noise''.
\end{issues}

\begin{figure}[htbp]\centering
  \fig[0.50]{kmeans_fail.png}
  \caption{K-means on two elongated, interleaved structures. The natural
  groups are obvious to the eye and unavailable to a partition by nearest
  centre: the objective can only be reduced by cutting across the arms. This
  is the failure that motivates every method in the rest of the workbook.}
  \label{fig:kmeansfail}
\end{figure}

\begin{figure}[htbp]\centering
  \fig[0.90]{kmeans_storyboard.png}
  \caption{The full K-means story in one grid, as used in the lecture: the
  loop, an initialisation that goes wrong, the failure on non-convex shapes,
  and the diagnostic that the objective can be lowered without the answer
  getting better.}
  \label{fig:kmeansstory}
\end{figure}

\subsection{Is K-means useful for chemical tagging?}

Yes --- with the right expectations, and in the literature it has been used
twice in ways that matter here.

\citet{Hogg:16} clustered $\sim\!10^5$ APOGEE stars in a 15-dimensional
abundance space with K-means, with no positional information, and recovered
groups that are genuinely phase-space structures. That result defines what
abundance-space overdensities can be at $\sim\!0.04$ dex precision. And
\citet{GarciaDias:18} applied K-means to $153\,847$ normalised APOGEE
spectra for spectral classification rather than for tagging, and the classes it
returns do track the stellar parameters --- dwarfs separate from giants, bulge
from halo. But the same paper's conclusion is the opposite of tidy: flux space
has no obvious groups, a discrete classification does not organise parameter
space neatly, and the solution moves with the initialisation. That is the same
warning as the row-order result of \S\ref{sec:validation}, in a different space.
K-means is useful when the groups are blobs; the spectra are not blobs, and the
abundance space of \S\ref{sec:benchmark} is only partly one.

In \textbf{Section~\ref{sec:benchmark}} the same algorithm is run on the same
member matrix as the other arms --- the 1\,002-row matrix of
\textbf{Table~\ref{tab:clusters}} in this chapter's summary box, and the
982-star deduplicated population of \textbf{Table~\ref{tab:honest}} and
\textbf{Table~\ref{tab:literature}} --- so that the comparison between ``the
classical tool'' and the newer ones is measured rather than asserted.

\begin{summary}[K-MEANS ON OUR MEMBER MATRIX]
\texttt{article/scripts/kmeans\_baseline.py} loads the DR19 sample, takes the
same member-only matrix the other benchmark arms use ($1\,002$ stars, 16
abundances, 25 clusters, $\geq\!5$ members), runs K-means with $n_{\rm
init}=10$ and reports the paper's merit functions alongside the SSE curve. The
measured result is worth stating plainly, because it is not the one the gospel
predicts: at $K=25$, K-means reaches \textbf{homogeneity 0.449}
(completeness 0.493, v-measure 0.470, accuracy 0.296) --- higher than t-SNE's
0.292 and EVoC's 0.409 on the same stars, and within reach of UMAP's 0.547.
Two caveats travel with it. At $K=2$ the same matrix gives homogeneity 0.228
with completeness 0.954: the ``families'' reading, in which almost everything
lands in one group. And raising $K$ to 40 lifts homogeneity to 0.505 while
completeness falls to 0.474 --- over-splitting buys homogeneity, which is why
the paper's metric pair must always be quoted together.
\end{summary}

\begin{issues}[EXERCISES]
\begin{enumerate}
\item Show that the assignment and update steps each cannot increase $J$, and
  conclude that the algorithm terminates. Then construct a one-dimensional
  example (three points, two centres) where the global optimum and Lloyd's
  fixed point differ, and draw the two partitions.
\item Implement k-means++ seeding in ten lines of NumPy and compare the final
  $J$ over 50 random datasets against a uniform-random initialisation. How
  much of the improvement comes from the seeding rule, and how much from the
  multiple restarts?
\item K-means with $K=2$ on our member matrix asks a different question from
  $K=25$: not ``which cluster is this star in'' but ``which family''. Run
  both and compare the homogeneity. Which of the two numbers would you quote
  to support strong chemical tagging, and why is that the wrong number?
\item Replace the Euclidean distance in the assignment step with cosine
  distance, keeping the mean update. What happens to the objective's
  guarantee of monotone decrease, and does the algorithm still terminate?
\end{enumerate}
\end{issues}
